
Given the gap between the PEFT community's assumptions and Mamba's architecture, we designed experiment h-e1 as a \textbf{MUST\_WORK existence check} --- a protocol that gives projection-only LoRA every possible advantage and treats failure as an unambiguous signal rather than a tuning problem. The design philosophy is: if LoRA is going to work on Mamba for classification at all, it should work under the conditions we construct.

The key diagnostic insight that shaped the methodology: we measured \textbf{per-epoch accuracy alongside loss} at every training epoch. Loss curves can look plausible --- oscillating, non-zero, apparently ''doing something'' --- while accuracy reveals complete learning failure. Three consecutive identical accuracy values across epochs cannot be rationalized as slow convergence or unlucky initialization.

\subsubsection{3.1 Mamba Architecture: Selective State Space Scan and Gradient Flow}

Mamba-130m (model identifier: \texttt{state-spaces/mamba-130m-hf}) is a 24-layer causal language model with hidden dimension \texttt{d\textbackslash \_model = 768}, inner dimension \texttt{d\textbackslash \_inner = 1536} (expansion factor 2), state dimension \texttt{d\textbackslash \_state = 16}, and discrete-time rank \texttt{dt\textbackslash \_rank = 48}. Each layer contains a MambaBlock with the following \texttt{nn.Linear} projection layers: \texttt{in\textbackslash \_proj} (768 $\rightarrow$ 3072), \texttt{out\textbackslash \_proj} (1536 $\rightarrow$ 768), \texttt{x\textbackslash \_proj} (1536 $\rightarrow$ 112), and \texttt{dt\textbackslash \_proj} (48 $\rightarrow$ 1536).

The MambaBlock also contains a \texttt{conv1d} layer and, critically, the \textbf{selective SSM scan}: the computation \texttt{h\textbackslash \_t = Ā(\$\textbackslash Delta\$, A)\$\textbackslash cdot h\_\{t-1\}\$ + B̄(\$\textbackslash Delta\$, B)\$\textbackslash cdot x\_t\$}, \texttt{y\textbackslash \_t = \$C\textbackslash cdot h\_t\$}. This scan is implemented either as a custom CUDA parallel associative scan (\texttt{selective\textbackslash \_scan\textbackslash \_cuda}) when the Mamba CUDA extensions are installed, or as a sequential Python fallback using standard PyTorch operations. Our experimental environment used the sequential fallback, as confirmed by the runtime warning logged at the start of every experimental run.

\textbf{Implementation note on gradient flow:} Because our experiments used the sequential implementation, the SSM scan was executed via standard PyTorch autograd operations rather than a custom CUDA kernel's backward pass. This means gradient flow through the scan was mediated by autograd in the standard manner. The failure of projection-layer LoRA to produce classification learning in this environment suggests that the gradient barrier effect is not solely attributable to custom CUDA kernel backward-pass limitations, and may instead reflect the geometric properties of gradients propagating through a pretrained SSM scan for a classification objective that diverges sharply from the pretraining task.

The gradient path for a classification loss placed on the final hidden state runs: \texttt{classification\textbackslash \_head \$\textbackslash rightarrow\$ final hidden state \$\textbackslash rightarrow\$ (through scan and projection layers of all 24 layers) \$\textbackslash rightarrow\$ LoRA weight matrices}. The SSM scan sits \textit{between} the classification head and the LoRA targets in every layer.

The \texttt{conv1d} layer is excluded from LoRA targets because it is not \texttt{nn.Linear}. The \texttt{dt\textbackslash \_proj} layer was excluded from the primary configuration as consistent with the standard community configuration.

\subsubsection{3.2 LoRA Application to Mamba-130m}

We applied LoRA using the HuggingFace PEFT library (v0.9+) with the following configuration, matching the community reference (alxndrTL/mamba-peft, 2024):

\begin{table}[htbp]
\centering
\resizebox{\columnwidth}{!}{%
\begin{tabular}{lll}
\toprule
Hyperparameter & Value & Rationale \\
\midrule
Target modules & `in\_proj`, `out\_proj`, `x\_proj` & Standard Mamba-1 projection LoRA; community reference configuration \\
Rank (r) & 8 & Standard for classification fine-tuning; referenced in MambaPEFT \\
Alpha ($\alpha )$ & 16 & 2r; standard scaling factor \\
Dropout & 0.05 & Standard regularization \\
Trainable parameters & 1,489,920 (1.14\% of 130M) & Confirmed by PEFT `print\_trainable\_parameters()` \\
\bottomrule
\end{tabular}
}
\end{table}

The LoRA adapters are installed into all 24 MambaBlocks. The classification head is a single \texttt{nn.Linear(768, num\textbackslash \_labels)} layer, applied to the last-token hidden state from the final MambaBlock output. The head is randomly initialized; no special initialization is applied.

\subsubsection{3.3 GLUE Fine-Tuning Setup}

We fine-tune separately on each GLUE task using the following shared training protocol:

\begin{table}[htbp]
\centering
\resizebox{\columnwidth}{!}{%
\begin{tabular}{ll}
\toprule
Hyperparameter & Value \\
\midrule
Optimizer & AdamW \\
Learning rate & $3\times 10^{-4}$ \\
Weight decay & 0.01 \\
Batch size & 32 \\
Epochs & 3 \\
LR schedule & Linear warmup (6\%) + linear decay \\
Seed & 42 \\
Training samples & 4,000 (subsampled from train split) \\
Max sequence length & 128 tokens \\
Tokenizer & GPT-NeoX-20B tokenizer (50k vocab) \\
\bottomrule
\end{tabular}
}
\end{table}

All four GLUE tasks (SST-2, MNLI, QNLI, QQP) were evaluated in experiment v8, which completed successfully. An earlier partial run (v8 first attempt) was interrupted before per-epoch evaluation completed for SST-2 due to an OOM error in a subsequent experiment version (v9) that attempted to add \texttt{dt\textbackslash \_proj} to the LoRA target modules; v8 was re-run and completed all four tasks.

\textbf{MUST\_WORK gate criterion:} We defined a conservative pass criterion: SST-2 accuracy > 70\% after three epochs. The conservatism is deliberate. The community reference reports 90--92\% --- a 70\% gate gives LoRA substantial benefit of the doubt, allowing for differences in training setup that might cost up to 20 percentage points of performance.

\subsubsection{3.4 Diagnostic Methodology}

Our core diagnostic design choice was \textbf{per-epoch accuracy measurement}, not just loss monitoring. Accuracy is the authoritative criterion because it measures task performance directly. Three consecutive identical accuracy values across epochs eliminate all alternative explanations:
\begin{itemize}
\item Slow convergence would produce monotonically increasing (if slow) accuracy
\item Lucky initialization plateau would not survive three consecutive identical values
\item Seed sensitivity would produce variation across epochs
\item Evaluation bug (same checkpoint evaluated repeatedly) was ruled out by confirming that evaluation was run as a per-epoch callback and that LoRA weight values differ between epochs confirming distinct model states were evaluated
\end{itemize}

We evaluate accuracy on the full validation split at the end of each epoch: SST-2 validation (872 examples), MNLI matched validation (872 examples, subsampled from 9,815), QNLI validation (872 examples), QQP validation (872 examples). We also log cross-entropy loss per training batch at every 20 steps.

\textbf{Zero-shot evaluation as baseline:} Before any fine-tuning, we evaluate the base Mamba-130m model on each GLUE task using lm-evaluation-harness with multiple-choice prompting. Zero-shot baselines are established in the same experimental environment as fine-tuning experiments.

